\section{Testing and Continuous Integration}\label{sec:testing-ci}

The project uses separate CI workflows for the paper, compiler, server, and client through GitHub Actions
workflows~\cite{GitHub_Actions_2026}.
The paper workflow checks LaTeX formatting, spelling, linting, and PDF build.
The compiler workflow checks formatting, builds the native compiler, runs CTest with GoogleTest, publishes test reports,
and summarizes coverage through gcovr.
The server and client workflows run ESLint, TypeScript type checks, Prettier checks, Vitest tests with coverage, builds,
and Docker image packaging.

\begin{figure}[htbp]
  \centering
  \begin{tikzpicture}[
      job/.style={rounded corners=4pt, draw=motivoInk!60, very thick, minimum width=2.05cm, minimum height=0.82cm,
      align=center, font=\small\sffamily},
      arrow/.style={-Stealth, thick, draw=motivoInk!70}
    ]
    \node[job, fill=motivoBlue!14] (format) {format};
    \node[job, fill=motivoCyan!14, right=0.28cm of format] (lint) {lint};
    \node[job, fill=motivoGreen!14, right=0.28cm of lint] (types) {types};
    \node[job, fill=motivoAmber!20, right=0.28cm of types] (test) {tests};
    \node[job, fill=motivoRose!12, right=0.28cm of test] (build) {build};
    \node[job, fill=motivoViolet!13, right=0.28cm of build] (package) {package};

    \draw[arrow] (format) -- (lint);
    \draw[arrow] (lint) -- (types);
    \draw[arrow] (types) -- (test);
    \draw[arrow] (test) -- (build);
    \draw[arrow] (build) -- (package);
  \end{tikzpicture}
  \caption{Common CI progression used across the project, adapted per component.}\label{fig:ci-flow}
\end{figure}

The test suites cover different failure modes.
Compiler unit tests use GoogleTest and exercise note mapping, parsing, semantic rules, lowering expressions, cursor
behavior, pattern calls, voice parallelism, and MIDI writing~\cite{GoogleTest_2026}.
Golden tests check representative success and error fixtures.
Server and client tests use Vitest to cover request validation, controller behavior, compiler process integration,
diagnostics parsing, editor setup, MIDI parsing, playback utilities, piano-roll rendering, shared UI components, and IDE
workflow hooks~\cite{Vitest_2026}.
This distribution is appropriate for a full-stack compiler tool because failures can happen at several
boundaries: language semantics, binary serialization, process execution, HTTP transport, browser parsing, and
UI rendering.

To ensure the correctness of the compilation pipeline and the validity of the generated MIDI artifacts, a comprehensive
testing suite was implemented using the \textit{Google Test} framework.
The testing strategy follows a tiered approach, covering individual units, module boundaries, and end-to-end
integration.

\subsection{Unit Testing}\label{detail-subsec:unit-testing}

Unit tests focus on the smallest verifiable components of the compiler.
This is particularly important for the musical domain logic, where small errors in calculation can lead to significant
musical dissonance.
The test cases cover several error-prone responsibilities.

\noindent \textbf{Pitch Math.} Verifying that pitch literals (e.g., \texttt{A4}) map correctly to MIDI note
numbers (e.g.,
69) and that octave offsets are handled according to the standard.

\noindent \textbf{Diagnostic Reporting.} Ensuring that the \texttt{DiagnosticsEngine} correctly captures and formats
errors across different stages of the pipeline.

\noindent \textbf{Lexer Accuracy.} Stress-testing the Flex regex patterns with edge cases, such as multiple consecutive
accidentals or varying whitespace.

\subsection{Integration and Regression Testing}\label{detail-subsec:integration-and-regression-testing}

Integration tests verify that the different stages of the compiler work together seamlessly.
These tests utilize a ``Golden File'' methodology:

\noindent A set of representative Motivo source files (located in \texttt{tests/data/}) is compiled.

\noindent The resulting IR is checked for temporal consistency (e.g., ensuring no notes overlap unless explicitly
requested).

\noindent The final MIDI binary is compared against a previously validated ``golden'' version to detect
unexpected changes
in behavior.

Regression tests were instrumental during the refactoring process (e.g., when decomposing the monolithic lowerer).
By maintaining a suite of complex compositions---such as the \texttt{fur\_elise.motivo} example---the compiler's
output could be
verified for byte-level consistency after every architectural change.

\subsection{Integration in the Development Lifecycle}\label{detail-subsec:integration-in-the-development-lifecycle}

The testing suite is integrated into the build system via CMake's \texttt{CTest} module.
This allows for one-command verification of the entire system (\texttt{make test}).
By enforcing a high level of test coverage, particularly in the lowering and backend stages, the implementation provides
repeatable evidence that the generated compositions are mathematically and musically accurate.
This verification process supports the thesis by treating the creative tool as a repeatable software system.

\begin{table}[htbp]
  \centering
  \renewcommand{\arraystretch}{1.3}
  \begin{tabularx}{\textwidth}{>{\raggedright\arraybackslash}p{2.8cm}>{\raggedright\arraybackslash}p{3.2cm}X}
    \toprule
    \textbf{Subsystem} & \textbf{Primary tools} & \textbf{What the tests protect} \\
    \midrule
    Compiler & CTest and GoogleTest & Grammar behavior, semantic diagnostics, lowering, MIDI writer, golden fixtures. \\
    Server & Vitest and Supertest & Request validation, compile route behavior, fake compiler integration,
    diagnostics parsing. \\
    Client & Vitest and Testing Library & Editor wiring, MIDI parsing, piano-roll rendering, playback hooks,
    IDE components. \\
    Paper & latexmk, tex-fmt, lint, codespell & Formatting, spelling, LaTeX warnings, and reproducible PDF
    generation. \\
    \bottomrule
  \end{tabularx}
  \caption{Testing responsibilities by subsystem.}\label{tab:testing-matrix}
\end{table}

This engineering chapter completes the description of the implemented system.
The next chapter evaluates Motivo through local measurements and source-to-output comparisons.
